\documentclass{beamer}

\usepackage{fvextra}
\newcommand{\aj}{The Astronomical Journal}
\newcommand{\apj}{ApJ}
\newcommand{\apjl}{ApJL}
\newcommand{\apjs}{ApJS}
\newcommand{\aap}{A\&A}
\newcommand{\mnras}{MNRAS}
\newcommand{\nar}{NAR}
\newcommand{\araa}{Annu. Rev. Astron. Astrophys.}
% Tema base
\usetheme{Madrid}

% Definición de colores (verde oscuro formal)
\definecolor{verdeoscuro}{RGB}{0,70,50}
\definecolor{verdeclaro}{RGB}{120,180,150}

% Aplicar colores al tema
\setbeamercolor{structure}{fg=verdeoscuro}
\setbeamercolor{palette primary}{bg=verdeoscuro, fg=white}
\setbeamercolor{palette secondary}{bg=verdeoscuro!85, fg=white}
\setbeamercolor{palette tertiary}{bg=verdeoscuro!70, fg=white}
\setbeamercolor{palette quaternary}{bg=verdeoscuro!50, fg=white}

\setbeamercolor{title}{fg=white}
\setbeamercolor{frametitle}{bg=verdeoscuro, fg=white}

\setbeamercolor{block title}{bg=verdeoscuro, fg=white}
\setbeamercolor{block body}{bg=verdeclaro!20, fg=black}

% Quitar símbolos de navegación
\setbeamertemplate{navigation symbols}{}

% Datos de la presentación
\title{Informe de proyecto de maestría}
\author{Tubal Caín Navarro Ruiz}
\institute{IRyA - UNAM}
\date{Abril 2025}

\begin{document}

% Portada
\begin{frame}
    \titlepage
\end{frame}

% Índice

\begin{frame}{Contenido}
    \tableofcontents
\end{frame}

% Sección 1

% Sección 2
\section{Scripts adaptados}

\begin{frame}{yso\_models.py}
    \begin{block}{Script de modelos}
        Crea modelos a partir de sf3dmodels
    \end{block}

    \begin{itemize}
        \item Escrito orientado a funciones. 
        \item Actualmente con 3 modelos: Ulrich+Disc, Hamburguer y Hamburguer Piecewise
        \item Ulrich probado con parámetros razonables 
    \end{itemize}
    
\end{frame}

\begin{frame}{make\_line.py}
    \begin{block}{Script de imagenes}
        Usa radmc3d para transferencia radiativa y manipula para crear fits convolucionados
    \end{block}
    \begin{itemize}
        \item Funciona con parser$^*$
        \item Usa datos observacionales para adaptar el modelo. (Distancias, tamaños, linea observada, beam, etc.)
        \item Corre radmc3d para generar el cubo de linea
        \item Corre Spectralcube para convolucionar; actualiza headers
        \item El output son imágenes crudas, convolucionadas en Jy/px y Jy/beam
    \end{itemize}
\end{frame}

\begin{frame}{main.py}
    \begin{block}{Script maestro}
        Ejecuta todos los pasos de los modelos
    \end{block}
    \begin{itemize}
        \item Crea una carpeta por modelo
        \item Llama a yso\_models y make\_line dentro de su respectiva carpeta
        \item Posiblemente a futuro ejecute mcmc o mande llamar a un script que lo ejecute
    \end{itemize}
\end{frame}


\section{Trabajo futuro}

\begin{frame}{Trabajo próximo}
\begin{itemize}
    \item Introducir y probar parámetros razonables en modelos Hamburguer y Hamburguer\_Piecewise
    \item Introducir ruido en la imagen final (make\_line), con formato adecuado para comparar con datos observacionales
    \item ¿Manipular datos observacionales?
    \item Crear script que ejecute mcmc
\end{itemize}
\end{frame}

\begin{frame}{Problemas específicos}
\begin{itemize}
    \item En los parámetros probados la abundancia es aun relativamente alta ($7\times 10^{-7}$)
    \item A abundancias menores requiere $\dot{M} $ mayor; pero con el grid actual genera concentraciones con $T>2000$ (pocos puntos). ¿Puede extrapolarse la Fn. de partición para no generar error al correr los modelos aunque no sea preciso? ¿Indica problemas en la resolución del modelo o de la RT?
    \item En los ejemplos de mcmc (rlg10) se usan custom keywords dentro del archivo radmc.inp. 
\end{itemize}
\end{frame}

\section{Bibliografía}
\begin{frame}[allowframebreaks]{Lecturas y consultas}
    \nocite{*}
    \bibliographystyle{apalike} % o plain, unsrt, etc.
    \bibliography{lecturas}
\end{frame}

% Frame final
\begin{frame}
    \centering
    \Huge Gracias
\end{frame}

\end{document}